\section{Representasi Data pada Level Mesin}
\label{sec:representasi-data}

Prosesor tidak mengenal huruf, tanda minus, maupun angka desimal. Ia
hanya mengenal tegangan tinggi dan rendah, yang dimodelkan sebagai
bit bernilai 1 dan 0. Karena itu seluruh data --- bilangan, huruf,
warna, bahkan instruksi --- harus disandikan menjadi deretan bit.
Bagian ini membahas penyandian yang paling sering dipakai.

\subsection{Sistem Bilangan Biner, Heksadesimal, dan Desimal}

Satu byte terdiri atas delapan bit dan dapat menyatakan
\dec{256} nilai berbeda, yaitu \dec{0} sampai \dec{255}. Angka
heksadesimal dipakai sebagai singkatan karena satu digit
heksadesimal tepat mewakili empat bit (satu \eng{nibble}).

\begin{table}[H]
    \centering
    \caption{Kesetaraan empat bit, digit heksadesimal, dan desimal.}
    \label{tab:hex-digit}
    \small
    \begin{tabular}{@{}cccc@{}}
        \toprule
        \textbf{Biner} & \textbf{Heksadesimal} & \textbf{Desimal} &
        \textbf{Bobot posisi} \\
        \midrule
        \bin{0000} & \hex{0} & 0 & $2^0 = 1$ \\
        \addlinespace
        \bin{0001} & \hex{1} & 1 & $2^1 = 2$ \\
        \addlinespace
        \bin{0010} & \hex{2} & 2 & $2^2 = 4$ \\
        \addlinespace
        \bin{0100} & \hex{4} & 4 & $2^3 = 8$ \\
        \addlinespace
        \bin{1000} & \hex{8} & 8 & \\
        \addlinespace
        \bin{1010} & \hex{A} & 10 & \\
        \addlinespace
        \bin{1111} & \hex{F} & 15 & \\
        \bottomrule
    \end{tabular}
\end{table}

Konversi biner ke heksadesimal dilakukan dengan mengelompokkan bit
per empat dari kanan. Konversi ke desimal dilakukan dengan
menjumlahkan bobot posisi setiap bit yang bernilai 1. Contoh:

\begin{align*}
    \bin{1011\,0110} &= \hex{B6} \\
    &= 11 \times 16 + 6 = 176 + 6 = 182
\end{align*}

Sebaliknya, konversi desimal ke heksadesimal dilakukan dengan
pembagian berulang oleh \dec{16} dan mencatat sisanya dari bawah ke
atas.

\subsection{Bilangan Bertanda dan \eng{Two's Complement}}

Bilangan biner di atas belum memiliki tanda. Untuk menyatakan
bilangan negatif, 8086 memakai representasi \term{two's complement}.
Aturannya sederhana: bilangan paling kiri (bit paling bermakna)
menjadi bit tanda. Nilai 0 berarti positif, nilai 1 berarti negatif.

\begin{definisi}[Two's complement]
Untuk mencari bentuk \eng{two's complement} dari sebuah bilangan,
balikkan seluruh bitnya (satu menjadi nol dan sebaliknya), lalu
tambahkan satu.
\end{definisi}

\begin{table}[H]
    \centering
    \caption{Representasi 8-bit: tak bertanda dan bertanda.}
    \label{tab:twos-complement}
    \small
    \begin{tabular}{@{}cccc@{}}
        \toprule
        \textbf{Biner} & \textbf{Tak bertanda} &
        \textbf{Bertanda (two's complement)} & \textbf{Keterangan} \\
        \midrule
        \bin{0000\,0000} & 0 & 0 & nol \\
        \addlinespace
        \bin{0111\,1111} & 127 & +127 & maksimum positif \\
        \addlinespace
        \bin{1000\,0000} & 128 & $-128$ & minimum negatif \\
        \addlinespace
        \bin{1111\,1111} & 255 & $-1$ & komplemen dari 1 \\
        \addlinespace
        \bin{1111\,0110} & 246 & $-10$ & komplemen dari 10 \\
        \bottomrule
    \end{tabular}
\end{table}

Rentang nilai yang dapat dinyatakan berbeda antara kedua tafsiran:
tak bertanda \dec{0} sampai \dec{255}, bertanda $-128$ sampai
$+127$. \important{Ketika menjumlahkan dua bilangan bertanda,
hasilnya bisa \emph{melampaui} rentang yang tersedia. Peristiwa ini
disebut \eng{overflow} dan ditandai oleh flag \flag{OF}.} Perhatikan
bahwa \eng{overflow} berbeda dari \eng{carry}: \eng{carry}
(\flag{CF}) menandakan kelebihan pada penjumlahan \emph{tak
bertanda}, sedangkan \eng{overflow} (\flag{OF}) menandakan kelebihan
pada penjumlahan \emph{bertanda}. Pembedaan ini dibahas tuntas pada
Bab~\ref{chap:register}.

\subsection{BCD dan ASCII}

Dua penyandian lain perlu dikenal karena erat kaitannya dengan
operasi masukan/keluaran pada 8086.

\begin{description}
    \item[BCD (\eng{Binary Coded Decimal})] Setiap digit desimal
          disandikan dalam empat bit. Angka \dec{47} menjadi
          \bin{0100\,0111}, bukan \bin{0010\,1111}. Penyandian ini
          memudahkan pengubahan ke tampilan desimal, dan operasi
          koreksinya disediakan oleh instruksi \asm{DAA} dan
          \asm{DAS}.

    \item[ASCII] Setiap karakter disandikan dalam tujuh atau delapan
          bit. Huruf \code{A} adalah \hex{41}, huruf \code{a} adalah
          \hex{61}, dan digit \code{7} adalah \hex{37}. Perhatikan
          bahwa digit karakter tidak sama dengan nilai bilangannya;
          konversi antara keduanya memerlukan penambahan atau
          pengurangan \hex{30}.
\end{description}

\begin{infobox}
Perbedaan antara nilai dan sandian karakter adalah sumber kesalahan
yang sangat sering terjadi pada pemrograman assembly. Menulis
\code{'7'} berarti menyimpan \hex{37}, sedangkan menulis \dec{7}
berarti menyimpan \hex{07}. Untuk mencetak sebuah bilangan,
program harus lebih dahulu mengubah nilai biner menjadi deretan
sandian karakter. Algoritmanya dibahas pada
Bab~\ref{chap:aritmatika}.
\end{infobox}

\subsection{Peta Konsep Buku}

Gambar~\ref{fig:peta-konsep} memperlihatkan bagaimana bab-bab buku
ini saling bergantung. Bab~III--VI membangun pemahaman tentang
mesin, Bab~VII--IX menjelaskan bagaimana instruksi berhubungan
dengan perangkat keras, Bab~X--XIII melatih penulisan program,
Bab~XIV melengkapi dengan operasi masukan/keluaran, dan
Bab~XV--XVI menutup dengan kemampuan analisis serta perbaikan.

\begin{figure}[H]
    \centering
    \begin{tikzpicture}[
        klaster/.style={
            draw=obeBlue, line width=0.9pt, rounded corners=4pt,
            inner sep=7pt, align=center, text width=3.7cm
        },
        hub/.style={-{Stealth[length=2.5mm]}, line width=1pt, obeTeal,
                    dashed}
    ]
        \node[klaster, fill=obeBlueLight] (k1) {
            {\small\bfseries CPMK01 --- Mesin}\\[0.3em]
            {\scriptsize Bab III--VI\\arsitektur, register,\\segmentasi memori}
        };
        \node[klaster, fill=obeAmberLight, draw=obeAmber,
              right=0.5cm of k1] (k2) {
            {\small\bfseries CPMK02 --- Instruksi}\\[0.3em]
            {\scriptsize Bab VII--IX, XIV\\format instruksi, mode\\pengalamatan, interupsi}
        };
        \node[klaster, fill=obeTealLight, draw=obeTeal,
              below=0.6cm of k1] (k3) {
            {\small\bfseries CPMK03 --- Menulis}\\[0.3em]
            {\scriptsize Bab X--XIII\\aritmatika, logika,\\kontrol, prosedur}
        };
        \node[klaster, fill=obeGreenLight, draw=obeGreen,
              right=0.5cm of k3] (k4) {
            {\small\bfseries CPMK04 --- Perbaikan}\\[0.3em]
            {\scriptsize Bab XV--XVI\\debugging, disassembly,\\optimasi}
        };

        \draw[hub] (k1) -- (k2);
        \draw[hub] (k1) -- (k3);
        \draw[hub] (k3) -- (k4);
        \draw[hub] (k3) -- (k2);
    \end{tikzpicture}
    \caption{Peta konsep buku: empat klaster kompetensi dan bab-bab
    yang menopangnya. Garis putus-putus menunjukkan ketergantungan
    materi.}
    \label{fig:peta-konsep}
\end{figure}
